\documentclass[a4paper,abstract,german]{scrreprt}

%
% Packages
%
\usepackage[german]{babel}
\usepackage[utf8]{inputenc}
\usepackage[T1]{fontenc}
\usepackage{lmodern}
\usepackage{amsmath}
\usepackage{bm}
\usepackage{amsmath}
\usepackage{amsfonts} 
\usepackage{amssymb}
\usepackage{amsthm}
\usepackage{graphicx}
\usepackage{enumitem}
\usepackage{tikz}
\usepackage[left=2.5cm,right=2.5cm,top=3cm,bottom=2cm,includeheadfoot]{geometry}

\setlength{\parindent}{0pt} 



%
% Eigene Macros
%
\newcommand{\N}{\mathbb{N}} 
\newcommand{\Z}{\mathbb{Z}} 
\newcommand{\Q}{\mathbb{Q}} 
\newcommand{\R}{\mathbb{R}} 
\newcommand{\C}{\mathbb{C}} 

%
% Beginn des eigentlichen Dokuments
%
\begin{document}


\noindent
{\begin{flushright}
%
% Hier die Namen einfüllen
%	
1. Andre Hemberger \\
2. Ferdinand Schneider \\
3. Anton Kohrt \\
4. Felix Krug \\
5. Luis Schreck \\
6. Jonathan Tschanter
\end{flushright}}
\begin{center}
	{\textbf  {Mathematik für Informatiker 2}} \\
	{\textbf {Abgabe 04. Übungsblatt}}
\end{center}\vspace{0.3cm}

%
% Aufzählung der Aufgaben
%
\begin{enumerate}

	\item[{1.}]
    \textbf{(2 Punkte)}
    Bestimmen Sie den links- und rechtsseitigen Grenzwert an der Stelle \( x_{0} = 2 \) von
    \[ 
        f(x) = \begin{cases}
            \hspace{0.2cm} x^2 + 1 & \text{ } x \geq 2 \\
            \hspace{0.2cm} x^2 + 1 & \text{ } x < 2.
        \end{cases} 
    \]
    Ist \textit{f} an der Stelle \(x_{0} = 2\) stetig? Skizzieren Sie den Graph von \textit{f}.
	

\end{enumerate}

\end{document}
